%------------------------------------------------------------------------------%

\usepackage{apresentacao}
\usepackage{mathconfig}
\usepackage{tabto}

%\setlength{\parskip}{0.5\baselineskip}

\newcounter{exemplo}
\setcounter{exemplo}{0}

\makeatletter
\graphicspath  {{./figs/}}
\def\input@path{{./figs/}}
\makeatother

\newcommand{\D}[2]{\dfrac{\partial {#1}}{\partial {#2}}}
\newcommand{\DS}[2]{\frac{\partial^2 {#1}}{\partial {#2}^2}}
\newcommand{\DM}[3]{\frac{\partial^2 {#1}}{\partial {#2} \partial {#3}}}

\newcommand{\vetor}[2]{\left(\begin{array}{c} #1 \\ #2 \end{array}\right)}

\newcommand{\veci}{\bm{i}}
\newcommand{\vecj}{\bm{j}}
\newcommand{\veck}{\bm{k}}

%------------------------------------------------------------------------------%
\title {Linearização}
\author{\large Luis Alberto D'Afonseca}
\date  {Cálculo de Funções de Várias Variáveis I}

%------------------------------------------------------------------------------%
\begin{document}

%------------------------------------------------------------------------------%
\begin{frame}
    \titlepage
\end{frame}

%------------------------------------------------------------------------------%
\section{Variação em uma Direção}

%------------------------------------------------------------------------------%
%\framedgraphic{Derivada Direcional}{derivada_direcional}{}

%------------------------------------------------------------------------------%
\begin{frame}{Variação em uma Dimensão}
  Em uma dimensão podemos estimar o valor da função em um ponto próximo à $x_0$
  usando a derivada em $x_0$

  \pause
  \[
    df = f'(x_0) ds
  \]

  \[
    \pause
    ds = x - x_0
  \]

  \[
    \pause
    df \approx f(x) - f(x_0)
  \]
\end{frame}

%------------------------------------------------------------------------------%
\begin{frame}{Variação em Várias Dimensões}
  Em várias dimensões usamos a derivada direcional
  \pause
  \[
    df
    = \left(D_u f\evalat{(x_0, y_0)}{}\right) ds
    \pause
    = \left(\nabla f\evalat{(x_0, y_0)}{} \cdot u\right) ds
  \]
  \pause
  \[
    ds = \sqrt{(x - x_0)^2 + (y-y_0)^2}
  \]
  \pause
  \[
    df \approx f(x, y) - f(x_0, y_0)
  \]
\end{frame}

%------------------------------------------------------------------------------%
\addtocounter{exemplo}{1}
\begin{frame}{Exemplo \theexemplo}
  Calcule quanto
  \[
    f(x, y, z) = y\sin(x) + 2yz
  \]
  irá variar se nos movermos 0,1 unidades a partir de \(P_0(0, 1, 0)\) \\
  na direção do ponto \(P_1(2, 2, -2)\)
\end{frame}

\begin{frame}{Exemplo \theexemplo}
  Direção do movimento
  \[
    \pause
    v
    \pause
    = \overrightarrow{P_0P_1}
    \pause
    = \left(\begin{array}{r} 2 \\ 2 \\ -2 \end{array}\right)
    - \left(\begin{array}{r} 0 \\ 1 \\  0 \end{array}\right)
    \pause
    = \left(\begin{array}{r} 2 \\ 1 \\  -2 \end{array}\right)
    \pause
    = 2\veci + \vecj -2\veck
  \]
  \pause
  Normalizando o vetor
  \[
    \pause
    u
    \pause
    = \frac{v}{\abs{v}}
    \pause
    = \frac{2\veci + \vecj -2\veck}{\sqrt{2^2 + 1^2 + (-2)^2}}
    \pause
    = \frac{2\veci + \vecj -2\veck}{\sqrt{9}}
    \pause
    = \frac{2}{3} \veci + \frac{1}{3} \vecj - \frac{2}{3} \veck
  \]
\end{frame}

\begin{frame}{Exemplo \theexemplo}
  \onslide<+->{Derivadas parciais}
  \begin{align*}
    \onslide<+->{\D{f}{x} & = \D{}{x}\left(y\sin(x) + 2yz\right)}
    \onslide<+->{ = y\cos(x) \\[5mm]}
    \onslide<+->{\D{f}{y} & = \D{}{y}\left(y\sin(x) + 2yz\right)}
    \onslide<+->{ = \sin(x) + 2z\\[5mm]}
    \onslide<+->{\D{f}{z} & = \D{}{z}\left(y\sin(x) + 2yz\right)}
    \onslide<+->{ = 2y}
  \end{align*}
\end{frame}

\begin{frame}{Exemplo \theexemplo}
  \onslide<+->{Gradiente de $f$}

  \[
    \onslide<+->{\nabla f}
    \onslide<+->{= (y\cos(x))\veci + (\sin(x) + 2z)\vecj + (2y)\veck}
    \onslide<+->{= \left(\begin{array}{c}
       y\cos(x) \\
       \sin(x) + 2z \\
       2y
    \end{array}\right)}
  \]

  \[
    \onslide<+->{\nabla f(0, 1, 0)}
    \onslide<+->{= (1\cos(0))\veci + (\sin(0) + 2\times 0)\vecj + (2\times 1)\veck}
    \onslide<+->{= \veci + 2\veck}
    \onslide<+->{= \left(\begin{array}{r} 1 \\ 0 \\ 2 \end{array}\right)}
  \]
\end{frame}

\begin{frame}{Exemplo \theexemplo}
  Derivada direcional
  \[
    \pause
    D_u f\evalat{(0, 1, 0)}{}
    \pause
    = \left(\nabla f\evalat{(0, 1, 0)}{} \cdot u\right)
    \pause
    = \left(\veci + 2\veck\right)
    \cdot
    \left(\frac{2}{3} \veci + \frac{1}{3} \vecj - \frac{2}{3} \veck\right)
  \]
  \[
    \hphantom{D_u f\evalat{(0, 1, 0)}{}}
    \pause
    = 1\times\frac{2}{3} + 0\times\frac{1}{3} + 2 \times\left(- \frac{2}{3}\right)
    \pause
    = \frac{2}{3} - \frac{4}{3}
    \pause
    = -\frac{2}{3}
  \]
\end{frame}

\begin{frame}{Exemplo \theexemplo}
  Variação
  \[
    \pause
    df
    \pause
    = \left(D_u f\evalat{(0, 1, 0)}{}\right) ds
    \pause
    = -\frac{2}{3}\times 0,1
    \pause
    = -\frac{2}{3} \frac{1}{10}
    \pause
    = -\frac{2}{30}
    \pause
    = -\frac{1}{15}
    \pause
    \approx -0,067
  \]
\end{frame}

%------------------------------------------------------------------------------%
\section{Aproximação Linear}

%------------------------------------------------------------------------------%
\begin{frame}{Teorema do Incremento}
  Suponha que as derivadas parciais de primeira ordem de $f(x, y)$ sejam
  definidas em uma região aberta $R$, que contem o ponto $(x_0, y_0)$
  \pause
  e que $f_x$ e $f_y$ sejam contínuas em $(x_0, y_0)$,
  \pause
  então a variação
  \[
    \pause
    \Delta z
    = f\left(x_0+\Delta x, y_0 + \Delta y\right)
    - f(x_0, y_0)
  \]
  \pause
  com \(\left(x_0+\Delta x, y_0 + \Delta y\right)\in R\),
  \pause
  satisfaz
  \[
    \Delta z
    = f_x(x_0, y_0) \Delta x
    + f_y(x_0, y_0) \Delta y
    + \varepsilon_1 \Delta x
    + \varepsilon_2 \Delta y
  \]
  \[
    \pause
    \varepsilon_1, \varepsilon_2 \to 0 \text{ quando } \Delta x, \Delta y \to 0
  \]
\end{frame}

%------------------------------------------------------------------------------%
\begin{frame}{Aproximação Linear}
  Aproximação linear de uma função diferenciável de duas variáveis
  em torno do ponto \((x_0, y_0)\)
  \pause
  \[
    L(x, y)
    \pause
    = f(x_0, y_0)
    + f_x(x_0, y_0)(x-x_0)
    + f_y(x_0, y_0)(y-y_0)
  \]

  \pause
  Aproximação linear padrão de $f$ em \((x_0, y_0)\)
  \[
    f(x, y) \approx L(x, y)
  \]
\end{frame}

%------------------------------------------------------------------------------%
\addtocounter{exemplo}{1}
\begin{frame}{Exemplo \theexemplo}
  Encontre a linearização de
  \[
    f(x, y) = x^2 - xy + \frac{y^2}{2} + 3
  \]
  no ponto \((3, 2)\)
\end{frame}

\begin{frame}{Exemplo \theexemplo}
  \onslide<+->{Avaliando as derivadas parciais}
  \begin{align*}
    \onslide<+->{f_x(x, y)}
    \onslide<+->{& = \D{}{x} \left(x^2 - xy + \frac{y^2}{2} + 3\right)}
    \onslide<+->{= 2x - y \\[5mm]}
    \onslide<+->{f_y(x, y)}
    \onslide<+->{& = \D{}{y} \left(x^2 - xy + \frac{y^2}{2} + 3\right)}
    \onslide<+->{= -x + y}
  \end{align*}
\end{frame}

\begin{frame}{Exemplo \theexemplo}
\onslide<+->{Avaliando a função e as derivadas parciais no ponto \((3, 2)\)}
\begin{align*}
  \onslide<+->{f(3, 2)}
  \onslide<+->{& = \left(x^2 - xy + \frac{y^2}{2} + 3\right)\evalat{(3, 2)}{}}
  \onslide<+->{= 3^2 - 3\times 2 + \frac{2^2}{2} + 3}
  \onslide<+->{= 8 \\[3mm]}
  \onslide<+->{f_x(3, 2)}
  \onslide<+->{& = \left(2x - y\right)\evalat{(3, 2)}{}}
  \onslide<+->{= 2\times 3 - 2}
  \onslide<+->{= 4\\[3mm]}
  \onslide<+->{f_y(3, 2)}
  \onslide<+->{& = \left(-x + y\right)\evalat{(3, 2)}{}}
  \onslide<+->{= -3 + 2}
  \onslide<+->{= -1}
\end{align*}
\end{frame}

\begin{frame}{Exemplo \theexemplo}
  \onslide<+->{Aproximação linear}
  \begin{align*}
    \onslide<+->{L(x, y)}
    \onslide<+->{& = f(x_0, y_0) + f_x(x_0, y_0)(x-x_0) + f_y(x_0, y_0)(y-y_0) \\[3mm]}
    \onslide<+->{& = f(3, 2) + f_x(3, 2)(x-3) + f_y(3, 2)(y-2) \\[3mm]}
    \onslide<+->{& = 8 + 4(x-3) + (-1)(y-2) \\[3mm]}
    \onslide<+->{& = 4x -y -2}
  \end{align*}
\end{frame}

%------------------------------------------------------------------------------%
\section{Erro na Aproximação Linear}

%------------------------------------------------------------------------------%
\begin{frame}{Erro na Aproximação Linear}
  Se $f$ possui primeira e segunda derivadas parciais
  \pause
  sobre um conjunto aberto
  contendo umm retângulo $R$ centrado em \((x_0, y_0)\)
  \pause
  e se $M$ é qualquer
  limitante superior para os valores de \(\abs{f_{xx}}\),  \(\abs{f_{yy}}\)
  e \(\abs{f_{xy}}\) em $R$
  \pause
  \[
    \abs{f_{xx}} \leq M \qquad
    \abs{f_{yy}} \leq M \qquad
    \abs{f_{xy}} \leq M \qquad
    \forall (x, y) \in R
  \]
  \pause
  então o erro $E(x, y)$ na linearização satisfaz
  \[
    \abs{E(x, y)} \leq \frac{M}{2} \big(\abs{x-x_0}+\abs{y-y_0}\big)^2
  \]
\end{frame}

%------------------------------------------------------------------------------%
\addtocounter{exemplo}{1}
\begin{frame}{Exemplo \theexemplo}
  Encontre um limitante superior para o erro na aproximação
  \[
    f(x, y) \approx L(x, y) = 4x -y -2
  \]
  no retângulo
  \[
    R: \abs{x-3} \leq 0,1 \,,\quad \abs{y-2} \leq 0,1
  \]
\end{frame}

\begin{frame}{Exemplo \theexemplo}
  Calculando os módulos das derivadas segundas
  \begin{align*}
    \onslide<+->{f_{xx}(x, y)}
    \onslide<+->{& = \DS{}{x} \left(x^2 - xy + \frac{y^2}{2} + 3\right) \\[3mm]}
    \onslide<+->{& = \D{}{x}\left(2x - y\right) \\[3mm]}
    \onslide<+->{& = 2 \\[3mm]}
    \onslide<+->{\abs{f_{xx}(x, y)}}
    \onslide<+->{& = 2}
  \end{align*}
\end{frame}

\begin{frame}{Exemplo \theexemplo}
  \begin{align*}
    \onslide<+->{f_{yy}(x, y)}
    \onslide<+->{& = \DS{}{y} \left(x^2 - xy + \frac{y^2}{2} + 3\right) \\[3mm]}
    \onslide<+->{& = \D{}{y}\left(-x + y\right) \\[3mm]}
    \onslide<+->{& = 1 \\[3mm]}
    \onslide<+->{\abs{f_{yy}(x, y)}}
    \onslide<+->{& = 1}
  \end{align*}
\end{frame}

\begin{frame}{Exemplo \theexemplo}
  \begin{align*}
    \onslide<+->{f_{xy}(x, y)}
    \onslide<+->{& = \DM{}{x}{y} \left(x^2 - xy + \frac{y^2}{2} + 3\right) \\[3mm]}
    \onslide<+->{& = \D{}{x}\left(-x + y\right) \\[3mm]}
    \onslide<+->{& = -1 \\[3mm]}
    \onslide<+->{\abs{f_{xy}(x, y)}}
    \onslide<+->{& = 1}
  \end{align*}
\end{frame}

\begin{frame}{Exemplo \theexemplo}
  \onslide<+->{Como o máximo das deriavas segundas em $R$ é 2}

  \onslide<+->{podemos escolher $M=2$, portanto}
  \begin{align*}
      \onslide<+->{\abs{E(x, y)}}
      \onslide<+->{& \leq \frac{M}{2} \big(\abs{x-x_0}+\abs{y-y_0}\big)^2 \\[3mm]}
      \onslide<+->{& = \frac{2}{2} \big(\abs{x-3}+\abs{y-2}\big)^2 \\[3mm]}
      \onslide<+->{& = \big(\abs{x-3}+\abs{y-2}\big)^2}
  \end{align*}
\end{frame}

\begin{frame}{Exemplo \theexemplo}
  \onslide<+->{Como
  \[
    \abs{x-3} \leq 0,1 \qquad \abs{y-2} \leq 0,1
  \]}
  \begin{align*}
    \onslide<+->{\abs{E(x, y)}}
    \onslide<+->{& \leq \big(\abs{x-3}+\abs{y-2}\big)^2\\[3mm]}
    \onslide<+->{& \leq \big(0,1+0,1\big)^2\\[3mm]}
    \onslide<+->{& = \big(0,2\big)^2\\[2mm]}
    \onslide<+->{& = 0,04}
  \end{align*}
\end{frame}

%------------------------------------------------------------------------------%
\section{Lista Mínima}

%------------------------------------------------------------------------------%
\begin{frame}{Lista Mínima}
  Cálculo Vol. 2 do Thomas 12\textsuperscript{a} ed. -- Seção 14.6
  \begin{enumerate}
    \item Estudar o texto da seção
    \item Resolver os exercícios: 20, 22, 24, 26, 30, 34, 36, 42, 48
  \end{enumerate}
  \vfill
  Atenção: A prova é baseada no livro, não nas apresentações
\end{frame}

%------------------------------------------------------------------------------%
\begin{frame}
  \tableofcontents
\end{frame}

%------------------------------------------------------------------------------%
\end{document}
%------------------------------------------------------------------------------%